\chapter{Experimentos y Resultados}\label{cap4:resultados}

Este capítulo presenta los resultados de los dos experimentos ejecutados en GPU
NVIDIA RTX~5050 (CUDA 12.6, PyTorch 2.4, SEED = 42). Los experimentos siguen
una complejidad creciente: validación analítica 1D (NB01) y validación 2D con
campo complejo (NB02).

% ──────────────────────────────────────────────────────────────────────────────
\section{NB01: Validación 1D - Helmholtz con solución analítica}
\label{sec:resultados1d}

\subsection{Configuración experimental}

El primer experimento valida la arquitectura SIREN frente a la solución analítica
exacta de la ecuación de Helmholtz 1D. Se empleó la arquitectura SIREN
$5 \times 64$ (${\sim}8{,}400$ parámetros) con los hiperparámetros de la
Tabla~\ref{tab:hiperparametros}. Los puntos de colocación son $N_c = 1{,}000$
puntos uniformes en $[0,1]$ y las condiciones de frontera Dirichlet son
$E(0) = 1$ y $E(1) = \cos(k)$.

\subsection{Resultados de precisión}

La Tabla~\ref{tab:resultados1d} presenta las métricas de precisión obtenidas
tras el entrenamiento completo Adam (15{,}000 épocas) seguido de L-BFGS
(202 iteraciones), con un tiempo total de 251 segundos.

\begin{table}[htbp]
\centering
\begin{threeparttable}
\caption{Métricas de precisión - PINN-SIREN 1D ($k = 2\pi$).}
\label{tab:resultados1d}
\begin{tabular}{lc}
\toprule
Métrica & Valor \\
\midrule
Error $L^2$ relativo (\%)    & \textbf{0.006} \\
MSE                          & $1.5 \times 10^{-9}$ \\
RMSE                         & $3.9 \times 10^{-5}$ \\
MAE                          & $3.50 \times 10^{-5}$ \\
Error máximo                 & $5.77 \times 10^{-5}$ \\
$R^2$                        & 1.000000 \\
Correlación de Pearson       & 1.000000 \\
\midrule
Épocas Adam                  & 15{,}000 / 15{,}000 \\
Iteraciones L-BFGS           & 202 / 500 \\
Tiempo total (s)             & 251 \\
\bottomrule
\end{tabular}
\begin{tablenotes}\small
    \item Evaluado en malla uniforme de 1{,}000 puntos en $[0,1]$.
    Arquitectura SIREN: 5 capas $\times$ 64 neuronas, $\omega_0 = 1.0$.
    Solución analítica: $E_\text{exacta}(x) = \cos(kx)$.
    GPU: NVIDIA RTX 5050, SEED = 42.
\end{tablenotes}
\end{threeparttable}
\end{table}

\subsection{Análisis y comparativa con el estado del arte}

El error $L^2 = 0.006\%$ supera el umbral de la tesis ($< 5\%$) por casi tres
órdenes de magnitud y representa un factor de mejora de $\times 415$ respecto al
resultado de \textcite{schoder2024helmholtz}, que reporta un error del $2.49\%$
para Helmholtz 1D con PINNs de activación \textit{tanh}. Esta mejora drástica
es atribuible a la ventaja intrínseca de la activación sinusoidal SIREN para
representar funciones oscilatorias: la red aprende la solución $\cos(kx)$ de
forma casi exacta, con $R^2 = 1.000000$ y correlación de Pearson $= 1.000000$
en toda la malla de evaluación.

La Figura~\ref{fig:resultados1d} muestra la comparación entre la predicción de
la red y la solución analítica, así como la distribución del error puntual.

\begin{figure}[htbp]
    \centering
    \includegraphics[width=0.9\textwidth]{resultados_pinn_1d.png}
    \caption{Validación PINN-SIREN para Helmholtz 1D.
    Panel izquierdo: comparación entre la predicción de la red y la solución
    analítica $E_\text{exacta}(x) = \cos(kx)$ en 1{,}000 puntos uniformes en
    $[0,1]$. Panel derecho: error puntual $|\hat{E}(x) - E_\text{exacta}(x)|$.
    Error $L^2 = 0.006\%$, $R^2 = 1.000000$. GPU: NVIDIA RTX 5050, SEED = 42.}
    \label{fig:resultados1d}
\end{figure}

\begin{figure}[htbp]
    \centering
    \includegraphics[width=\textwidth]{metricas_adicionales_1d.png}
    \caption{Métricas adicionales de entrenamiento NB01. De izquierda a derecha:
    curva de pérdida total durante las fases Adam y L-BFGS, distribución del
    error puntual y métricas de ajuste ($R^2$, Pearson). La transición
    Adam $\to$ L-BFGS a la época 15{,}000 es visible como una caída abrupta
    en la curva de pérdida.}
    \label{fig:metricas1d}
\end{figure}

% ──────────────────────────────────────────────────────────────────────────────
\section{NB02: Validación 2D - Helmholtz con campo complejo}
\label{sec:resultados2d}

\subsection{Configuración experimental}

El segundo experimento escala el modelo al caso bidimensional con campo complejo.
Se empleó la arquitectura SIREN $5 \times 128$ ($66{,}690$ parámetros), LHS
con $N_c = 3{,}000$ puntos de colocación y $N_b = 300$ puntos de frontera por
borde. El peso de física $\lambda_\text{fís} = 0.1$ fue seleccionado mediante
el experimento de ablación descrito en la Sección~\ref{sec:ablacion}.

\subsection{Resultados de precisión}

La Tabla~\ref{tab:resultados2d} reporta las métricas detalladas para cada
componente del campo complejo, evaluadas en una malla uniforme de
$100 \times 100$ puntos.

\begin{table}[htbp]
\centering
\begin{threeparttable}
\caption{Métricas de precisión - PINN-SIREN 2D ($k = 2\pi$, $k_x = k_y = k/\sqrt{2}$).}
\label{tab:resultados2d}
\begin{tabular}{lcc}
\toprule
Métrica & $E_\text{real}$ & $E_\text{imag}$ \\
\midrule
Error $L^2$ (\%)         & 0.214              & 0.127              \\
MSE                      & $2.20 \times 10^{-6}$ & $8.43 \times 10^{-7}$ \\
RMSE                     & $1.48 \times 10^{-3}$ & $9.18 \times 10^{-4}$ \\
MAE                      & $1.27 \times 10^{-3}$ & $7.71 \times 10^{-4}$ \\
Error máximo             & $3.05 \times 10^{-3}$ & $2.93 \times 10^{-3}$ \\
$R^2$                    & 0.999995           & 0.999998           \\
Correlación de Pearson   & 0.999998           & 0.999999           \\
\midrule
Error $L^2$ promedio (\%)& \multicolumn{2}{c}{\textbf{0.171}} \\
\bottomrule
\end{tabular}
\begin{tablenotes}\small
    \item Evaluado en malla uniforme de $100 \times 100$ puntos.
    Arquitectura: SIREN $5 \times 128$, $\omega_0 = 1.0$, 66{,}690 parámetros.
    Puntos de colocación: LHS $N_c = 3{,}000$.
    Frontera: onda plana $e^{i(k_x x + k_y y)}$ en 4 bordes, $N_b = 300$ pts/borde.
    Entrenamiento: Adam (8{,}737 épocas) + L-BFGS (1{,}035 iter.), tiempo total 299 s.
    Varianza entre semillas (SEED $\in \{42, 123, 777\}$):
    error $L^2$ promedio $= 0.155 \pm 0.020$\,\%.
\end{tablenotes}
\end{threeparttable}
\end{table}

\subsection{Análisis de resultados}

El error $L^2$ promedio de $0.171\%$ supera el umbral de la tesis ($< 5\%$) por
más de un orden de magnitud y mejora en un factor de $\times 11.2$ el resultado
de \textcite{schoder2024helmholtz} para Helmholtz 2D ($1.91\%$). La precisión
se mantiene alta en ambas componentes del campo complejo, con $R^2 > 0.9999$
para $E_\text{real}$ y $E_\text{imag}$.

La ampliación de capacidad de $d = 64$ a $d = 128$ fue determinante: con
$d = 64$ el error promedio era $0.436\%$; con $d = 128$ se redujo a $0.171\%$,
una mejora de $2.5\times$. El Muestreo Latino Hipercúbico contribuyó a reducir
el error en regiones de baja densidad de puntos respecto a mallas cartesianas
equivalentes.

La Figura~\ref{fig:resultados2d} muestra los mapas de calor de la solución
exacta y la predicción PINN-SIREN para ambas componentes del campo.

\begin{figure}[htbp]
    \centering
    \includegraphics[width=0.9\textwidth]{resultados_pinn_2d.png}
    \caption{Validación PINN-SIREN para Helmholtz 2D con campo complejo.
    Fila superior: partes real e imaginaria de la solución analítica
    $E_\text{exacta} = e^{i(k_x x + k_y y)}$.
    Fila inferior: predicción PINN-SIREN correspondiente.
    La diferencia visual es indistinguible; los errores puntuales son del
    orden de $10^{-3}$ (véase la Tabla~\ref{tab:resultados2d}).
    Malla de evaluación $100 \times 100$, SEED = 42.}
    \label{fig:resultados2d}
\end{figure}

\begin{figure}[htbp]
    \centering
    \includegraphics[width=\textwidth]{metricas_adicionales_2d.png}
    \caption{Métricas adicionales de entrenamiento NB02. Panel izquierdo:
    curva de pérdida total mostrando la transición bifásica Adam $\to$ L-BFGS.
    Paneles central y derecho: mapas de error puntual
    $|E_\text{real} - \hat{E}_\text{real}|$ y
    $|E_\text{imag} - \hat{E}_\text{imag}|$ en la malla $100 \times 100$.
    Los errores más altos se concentran en las esquinas del dominio.}
    \label{fig:metricas2d}
\end{figure}

\subsection{Reproducibilidad multi-semilla}

Para evaluar la robustez estadística del método, el experimento NB02 se repitió
con tres semillas diferentes: SEED $\in \{42, 123, 777\}$. Los resultados se
presentan en la Tabla~\ref{tab:multiseed}.

\begin{table}[htbp]
\centering
\begin{threeparttable}
\caption{Reproducibilidad multi-semilla - NB02 ($\lambda_\text{fís} = 0.1$).}
\label{tab:multiseed}
\begin{tabular}{lccc}
\toprule
SEED & Error $L^2$ prom. (\%) & Épocas Adam & Tiempo (s) \\
\midrule
42   & 0.171  & 8{,}737  & 299 \\
123  & 0.1264 & 8{,}012  & 448 \\
777  & 0.1675 & 13{,}560 & 346 \\
\midrule
Media $\pm$ Desv. est. & $0.155 \pm 0.020$ & — & — \\
\bottomrule
\end{tabular}
\begin{tablenotes}\small
    \item Misma arquitectura, hiperparámetros y protocolo de entrenamiento.
    La varianza entre semillas es del 13\% relativo, confirmando la
    estabilidad del método. GPU: NVIDIA RTX 5050.
\end{tablenotes}
\end{threeparttable}
\end{table}

La varianza inter-semilla de $\pm 0.020\%$ (13\% relativo) confirma que el
método es estadísticamente robusto: independientemente de la inicialización
aleatoria, el modelo converge a soluciones de alta precisión.

% ──────────────────────────────────────────────────────────────────────────────
\section{Ablación del peso de física $\lambda_\text{fís}$}
\label{sec:ablacion}

Para calibrar el balance entre el término de datos y el término físico en la
función de pérdida~\eqref{eq:loss_2d}, se realizó un experimento de ablación
con $\lambda_\text{fís} \in \{0.01, 0.1, 1.0\}$, manteniendo fija la
arquitectura SIREN $5 \times 128$ y SEED = 42. Los resultados se presentan en
la Tabla~\ref{tab:ablacion_lambda}.

\begin{table}[htbp]
\centering
\begin{threeparttable}
\caption{Ablación del peso de física $\lambda_\text{fís}$ - NB02.}
\label{tab:ablacion_lambda}
\begin{tabular}{lccc}
\toprule
$\lambda_\text{fís}$ & Error $L^2$ prom. (\%) & Converge & Observación \\
\midrule
0.01  & 0.163  & Sí & Física subponderada \\
0.1   & 0.171  & Sí & \textbf{Configuración óptima} \\
1.0   & $>5.0$ & No & Física sobreponderada \\
\bottomrule
\end{tabular}
\begin{tablenotes}\small
    \item Arquitectura SIREN $5 \times 128$, $\omega_0 = 1.0$, SEED = 42.
    Mismo protocolo Adam + L-BFGS que NB02.
    \textit{No converge} indica que Adam no alcanzó la condición de parada
    antes de las 15{,}000 épocas con $\varepsilon_{L^2} < 5\%$. El caso
    $\lambda = 1.0$ produce una explosión numérica en la función de pérdida
    debido a que los dos residuos (real e imaginario) duplican la contribución
    de la física, generando gradientes de gran magnitud en \texttt{float32}.
\end{tablenotes}
\end{threeparttable}
\end{table}

El valor $\lambda_\text{fís} = 0.1$ equilibra adecuadamente el ajuste a los
datos de frontera y el cumplimiento de la EDP: el modelo con $\lambda = 0.01$
obtiene ligeramente mejor error (0.163\%) al sacrificar restricción física,
mientras que $\lambda = 1.0$ produce gradientes explosivos. Se seleccionó
$\lambda = 0.1$ como configuración óptima por su estabilidad y precisión.

% ──────────────────────────────────────────────────────────────────────────────
\section{Comparativa con el estado del arte}

La Tabla~\ref{tab:comparativa} resume el rendimiento de la PINN-SIREN propuesta
frente al estado del arte en la resolución de la ecuación de Helmholtz con PINNs.

\begin{table}[H]
\centering
\begin{threeparttable}
\caption{Comparativa de error $L^2$ relativo - PINN-SIREN vs.\ estado del arte.}
\label{tab:comparativa}
\begin{tabular}{lccc}
\toprule
Método & Dimensión & Error $L^2$ (\%) & Factor de mejora \\
\midrule
\textcite{schoder2024helmholtz} & 1D & 2.490 & -- \\
\textcite{schoder2024helmholtz} & 2D & 1.910 & -- \\
\textbf{PINN-SIREN (este trabajo)} & \textbf{1D} & \textbf{0.006} & $\times 415$ \\
\textbf{PINN-SIREN (este trabajo)} & \textbf{2D} & \textbf{0.171} & $\times 11.2$ \\
\bottomrule
\end{tabular}
\begin{tablenotes}\small
    \item El factor de mejora es la razón entre el error de Schoder et al.\ (2024)
    y el error de la PINN-SIREN propuesta.
    Ambos trabajos abordan la ecuación de Helmholtz escalar en dominios normalizados.
    \textcite{schoder2024helmholtz} aborda el caso 3D con activación \textit{tanh};
    este trabajo opera en 1D/2D adimensionales con arquitectura SIREN.
    El factor de mejora refleja tanto la ventaja arquitectural de SIREN sobre
    \textit{tanh}/ReLU como la diferencia de complejidad dimensional.
\end{tablenotes}
\end{threeparttable}
\end{table}

% ──────────────────────────────────────────────────────────────────────────────
\section{Resumen de experimentos completados}

La Tabla~\ref{tab:resumen_experimentos} sintetiza el estado de los experimentos
ejecutados en el marco de este primer avance.

\begin{table}[H]
\centering
\begin{threeparttable}
\caption{Estado de los experimentos ejecutados.}
\label{tab:resumen_experimentos}
\begin{tabular}{llcc}
\toprule
Experimento & Descripción & Estado & Métrica principal \\
\midrule
NB01 & Helmholtz 1D - validación analítica    & Completo  & $L^2 = 0.006\%$ \\
NB02 & Helmholtz 2D - campo complejo          & Completo  & $L^2 = 0.171\%$ \\
NB03 & Speckle óptico - frontera rugosa       & Pendiente & $|C - 1| < 0.1$ \\
NB04 & Benchmark FEM - factor de aceleración  & Pendiente & $S = T_\text{FEM}/T_\text{PINN}$ \\
\bottomrule
\end{tabular}
\begin{tablenotes}\small
    \item NB01 y NB02 verifican la hipótesis de precisión ($L^2 < 5\%$).
    NB03 y NB04 constituyen el trabajo futuro inmediato.
\end{tablenotes}
\end{threeparttable}
\end{table}



Los dos experimentos completados verifican la primera hipótesis de la tesis:
el error $L^2 < 5\%$ se supera con amplitud tanto en el caso 1D ($L^2 = 0.006\%$)
como en el 2D ($L^2 = 0.171\%$). La validación estadística del
\textit{speckle} y el \textit{benchmark} contra FEM se abordarán como
trabajo futuro inmediato.